% Part: incompleteness
% Chapter: incompleteness-provability
% Section: introduction

\documentclass[../../../include/open-logic-section]{subfiles}

\begin{document}

\olfileid{inc}{inp}{int}

\olsection{引言}

Hilbert 认为，一个数学结构（如自然数）的公理系统，如果不能推出关于该结构的所有真语句，就是不充分的。
结合他后来对形式演绎系统的兴趣，这表明他认为我们应当保证，例如，我们用来对自然数进行推理的形式系统不仅是一致的，而且是\emph{完备的}；
也就是说，其语言中的每个语句要么可推导，要么其否定可推导。
Gödel 第一不完全性定理表明，这样的公理系统不存在：
不存在对算术而言完备、一致且可公理化的形式系统。事实上，
任何“足够强的”、一致且可公理化的数学理论都不是完备的。

Hilbert 的一个更重要的目标，即其证立现代（“经典”）数学之纲领的核心，
是为代表经典推理的形式系统找到有穷一致性证明。
就 Hilbert 纲领而言，Gödel 的第二不完全性定理是一个严重得多的打击。
与第一不完全性定理一样，第二不完全性定理也可以用宽泛的措辞来陈述。
粗略地说，它指出没有足够强的算术理论能够证明自身的一致性。
我们必须将“足够强”理解为包含比~$\Th{Q}$ 稍多一点的内容。

Gödel 对不完全性定理最初证明背后的思想，可以在 Epimenides（埃庇米尼得斯） 悖论中找到。
Epimenides 是一个克里特人，他断言所有克里特人都是说谎者；
该悖论更直接的形式是断言“这句话是假的”。
本质上，通过将“真”替换为“可推导性”，
Gödel 得以形式化一个语句，该语句以迂回的方式断言其自身是不可推导的。
如果该语句是可推导的，那么该理论就会不一致。
Gödel 证明了该语句的否定也不能从他当时考虑的公理系统中推导出来。
（对于这第二部分，Gödel 必须假设理论~$\Th{T}$ 具有所谓的“$\omega$-一致性”。
$\omega$-一致性与一致性相关，但是一个更强的性质。\footnote{也就是说，任何 $\omega$-一致的理论都是一致的，但反之不然。}
几年后，Rosser 证明了只需假设~$\Th{T}$ 是简单一致的即可。）

第一个挑战是理解如何构造一个指涉自身的语句。
对于语言~$\Th{Q}$ 中的每个公式~$!A$，
令~$\gn{!A}$ 表示对应于~$\Gn{!A}$ 的数码。
想想这意味着什么：$!A$~是语言~$\Th{Q}$ 中的一个公式，
$\Gn{!A}$~是一个自然数，
而 $\gn{!A}$~是语言~$\Th{Q}$ 中的一个\emph{项}。
因此，语言~$\Th{Q}$ 中的每个公式~$!A$ 都有一个\emph{名字}~$\gn{!A}$，
它是语言~$\Th{Q}$ 中的一个项；
这为我们提供了一个概念框架，使得语言~$\Th{Q}$ 中的公式可以“谈论”其他公式。
下面的引理被称为不动点引理。

\begin{lem}
设 $\Th{T}$ 为任意扩张~$\Th{Q}$ 的理论，
并设 $!B(x)$ 为仅有变元~$x$ 自由的任意公式。
则存在语句~$!A$ 使得 $\Th{T} \Proves!A \liff !B(\gn{!A})$。
\end{lem}

该引理断言，给定任何性质 $!B(x)$，
总存在一个语句~$!A$ 断言“$!B(x)$ 对我成立”，
且 $\Th{T}$“知道”这一点。

我们如何构造这样的语句呢？考虑下面这个由 Quine 提出的
Epimenides 悖论变体：

\begin{quote}
“前面加上自身引文则得出假”前面加上自身引文则得出假。
\end{quote}

这个句子并非直接自指。它仅仅是对引号之间的句法对象做出了断言，
在这一点上，它与下列句子类似：

\begin{enumerate}
\item “Robert”是个好名字。
\item “我跑了。”是个短句。
\item “有三个词”有三个词。
\end{enumerate}

但是，如果我们把短语“前面加上自身引文则得出假”，
在它前面加上它自身的引文，会发生什么呢？
得到的恰好就是原句！简言之，这个句子断言它自身是假的。

\end{document}